\section{Transfer through relative interpretations}
\label{sec:annex-p05-interpretation-transfer}

\begin{annexitem}{annex:p05.interpretation-transfer}{thm:p05-interpretation-transfer}{Transfer}
\begin{proof}
Use the data of \cref{def:p05-relative-interpretation}. If \(M\models V\),
the domain formula defines a nonempty subset of \(M\). The interpreted
equality is an equivalence relation there. Take its equivalence classes
as the elements of \(\mathcal I(M)\). Invariance of the relation graphs
makes the interpreted relations well-defined on these classes.
Totality, uniqueness modulo interpreted equality, and invariance of
the function graphs give total functions on the classes, including
the nullary functions interpreting constants. These constructions use
the classes themselves, not a definable selection of representatives.

Induction on terms, using the defining graphs, and then on formulas
gives the truth equivalence
\[
 \mathcal I(M)\models\psi[\bar a/E]
 \quad\Longleftrightarrow\quad
 M\models\psi^{\mathcal I}[\bar a]
\]
for tuples \(\bar a\) in the interpreted domain. The equality case
is precisely the definition of the quotient. The other atomic cases
use graph invariance; Boolean cases use the commuting translation.
In a quantifier case, every equivalence class has a representative in
the domain, and the translated quantifier ranges over that domain.
Thus the quantified cases also agree. Only finitely many representatives
are needed for the free variables of an individual formula; no global
choice of representatives is assumed.

Since \(V\) proves each translated axiom of \(U\), soundness and the
truth equivalence show that \(\mathcal I(M)\models U\). If
\(U\derives\psi\), it follows that every \(M\models V\) satisfies
\(\psi^{\mathcal I}\). Strong first-order completeness therefore gives
\(V\derives\psi^{\mathcal I}\). If \(V\) is consistent, completeness
provides such an \(M\), and \(\mathcal I(M)\) is a model of \(U\).
Soundness then proves consistency of \(U\).

For the effective assertion, take an arbitrary consistent c.e. theory
\(V'\supseteq V\) in the same target signature. The interpretation
conditions remain provable in \(V'\). In the effective sentence codings
implicit in the recursive translation and c.e. theories, put
\[
 U' = \{\psi:\psi\text{ is a source sentence and }
                         V'\derives\psi^{\mathcal I}\}.
\]
This is a c.e. set. Enumerate source sentences and their computable
translations, and dovetail this with enumeration of \(V'\)-proofs.
For a c.e. axiom set, the latter enumeration may attach finite stages
of the axiom enumeration as certificates for axiom occurrences.
Every proof uses finitely many such occurrences and eventually appears.
Equality of a displayed conclusion and a computed translation is a
decidable finite-string test. This enumerates exactly \(U'\).

The first part of the proof gives \(U\subseteq U'\). A model of
\(V'\) exists by consistency and completeness; its interpreted model
satisfies every member of \(U'\), so \(U'\) is consistent. It is also
deductively closed. If finitely many sentences of \(U'\) entail a
sentence \(\theta\), then each interpreted model of a \(V'\)-model
satisfies those premises and hence \(\theta\). The truth equivalence
and completeness give \(V'\derives\theta^{\mathcal I}\), so
\(\theta\in U'\).

If theoremhood in \(V'\) were decidable, computing
\(\psi^{\mathcal I}\) and applying that decider would decide membership
in \(U'\). Deductive closure makes this also a decision of theoremhood
in the consistent c.e. extension \(U'\) of \(U\), contradicting the
essential undecidability of \(U\). Therefore every consistent c.e.
extension \(V'\) is undecidable. This is essential undecidability of
\(V\); because \(V\) itself is consistent and c.e., it is undecidable
as well. The use of essential undecidability is exact: \(U'\) need not
have the same theorems as \(U\), but is a consistent effective extension.
\end{proof}
\end{annexitem}
